\documentclass[10pt, a4paper]{report}
\usepackage[utf8]{inputenc}
\usepackage{geometry}
\usepackage{longtable}
\usepackage{array} % For better column definitions in longtable
\usepackage{enumitem} % Required for itemize and enumerate options
\usepackage{helvet} % Using a sans-serif font for better readability
\renewcommand{\familydefault}{\sfdefault}
\usepackage{xcolor}
\usepackage{fancyhdr}
\usepackage{hyperref}

\geometry{left=1.5cm, right=1.5cm, top=2cm, bottom=2cm}

\hypersetup{
    colorlinks=true,
    linkcolor=blue,
    filecolor=magenta,      
    urlcolor=cyan,
}

\pagestyle{fancy}
\fancyhf{}
\fancyhead[L]{Library Management System - Functional Test Cases}
\fancyfoot[C]{\thepage}

\title{Functional Test Cases \\ \large Library Management System}
\author{Test Design Team (Generated by Gemini)}
\date{\today}

\begin{document}
\maketitle
\begin{center}
    \textbf{Note:} These test cases are designed from a black-box, user-centric perspective, based on existing unit test files. Actual results and status are to be filled during test execution.
\end{center}

\section{Server-Side Functional Tests}

\subsection{Authentication \& Authorization}

\begin{longtable}{|>{\raggedright\arraybackslash}p{3.5cm}|>{\raggedright\arraybackslash}p{5cm}|>{\raggedright\arraybackslash}p{6cm}|>{\raggedright\arraybackslash}p{6cm}|}
\caption{Authentication and Authorization Test Cases}\label{tab:auth_test_cases}\\
\hline
\textbf{Test Case ID} & \textbf{Test Title} & \textbf{Test Steps} & \textbf{Expected Results} \\
\hline
\endfirsthead
\multicolumn{4}{c}%
{{\bfseries \tablename\ \thetable{} -- continued from previous page}} \\
\hline
\textbf{Test Case ID} & \textbf{Test Title} & \textbf{Test Steps} & \textbf{Expected Results} \\
\hline
\endhead
\hline \multicolumn{4}{|r|}{{Continued on next page}} \\ \hline
\endfoot
\hline
\endlastfoot

% User Registration
\multicolumn{4}{|l|}{\textbf{Sub-Section: User Registration}} \\ \hline
TC\_AUTH\_REG\_001 & Successful new user registration & 
    \begin{itemize}[noitemsep, topsep=0pt, partopsep=0pt, leftmargin=*]
        \item Send a request to the registration endpoint.
        \item Provide a valid and unique email, a password, and a name.
    \end{itemize} & 
    \begin{itemize}[noitemsep, topsep=0pt, partopsep=0pt, leftmargin=*]
        \item System accepts the data.
        \item A new user record is created in the database.
        \item The API returns the new user's data with a status of 201.
    \end{itemize} \\ \hline
TC\_AUTH\_REG\_002 & Attempt to register with an existing email & 
    \begin{itemize}[noitemsep, topsep=0pt, partopsep=0pt, leftmargin=*]
        \item Send a request to the registration endpoint with an email that already exists.
    \end{itemize} & 
    \begin{itemize}[noitemsep, topsep=0pt, partopsep=0pt, leftmargin=*]
        \item System rejects the registration.
        \item An error message "Email already exists" is returned.
        \item The API returns a status of 400.
    \end{itemize} \\ \hline
TC\_AUTH\_REG\_003 & Handle server error during registration &
    \begin{itemize}[noitemsep, topsep=0pt, partopsep=0pt, leftmargin=*]
        \item Trigger a server-side error during the user creation process (e.g., database connection fails).
    \end{itemize} &
    \begin{itemize}[noitemsep, topsep=0pt, partopsep=0pt, leftmargin=*]
        \item System handles the error gracefully.
        \item A generic error message is returned.
        \item The API returns a status of 400.
    \end{itemize} \\ \hline

% User Login
\multicolumn{4}{|l|}{\textbf{Sub-Section: User Login}} \\ \hline
TC\_AUTH\_LOGIN\_001 & Successful login with valid credentials & 
    \begin{itemize}[noitemsep, topsep=0pt, partopsep=0pt, leftmargin=*]
        \item Send a request to the login endpoint with a valid email and correct password.
    \end{itemize} & 
    \begin{itemize}[noitemsep, topsep=0pt, partopsep=0pt, leftmargin=*]
        \item User is authenticated successfully.
        \item A JWT token is returned in the response.
        \item The API returns a status of 200.
    \end{itemize} \\ \hline
TC\_AUTH\_LOGIN\_002 & Attempt login with invalid password & 
    \begin{itemize}[noitemsep, topsep=0pt, partopsep=0pt, leftmargin=*]
        \item Send a request to the login endpoint with a valid email but an incorrect password.
    \end{itemize} & 
    \begin{itemize}[noitemsep, topsep=0pt, partopsep=0pt, leftmargin=*]
        \item Login fails.
        \item An error message "Incorrect password" is returned.
        \item The API returns a status of 401.
    \end{itemize} \\ \hline
TC\_AUTH\_LOGIN\_003 & Attempt login with non-existent email & 
    \begin{itemize}[noitemsep, topsep=0pt, partopsep=0pt, leftmargin=*]
        \item Send a request to the login endpoint with an email that does not exist in the database.
    \end{itemize} & 
    \begin{itemize}[noitemsep, topsep=0pt, partopsep=0pt, leftmargin=*]
        \item Login fails.
        \item An error message "Cannot find user" is returned.
        \item The API returns a status of 400.
    \end{itemize} \\ \hline
TC\_AUTH\_LOGIN\_004 & Handle server error during login &
    \begin{itemize}[noitemsep, topsep=0pt, partopsep=0pt, leftmargin=*]
        \item Trigger a server-side error during the login process (e.g., database lookup fails).
    \end{itemize} &
    \begin{itemize}[noitemsep, topsep=0pt, partopsep=0pt, leftmargin=*]
        \item System handles the error gracefully.
        \item A generic error message is returned.
        \item The API returns a status of 500.
    \end{itemize} \\ \hline
\end{longtable}

\newpage
\subsection{Author Management}

\begin{longtable}{|>{\raggedright\arraybackslash}p{4cm}|>{\raggedright\arraybackslash}p{5cm}|>{\raggedright\arraybackslash}p{6cm}|>{\raggedright\arraybackslash}p{6cm}|}
\caption{Author Management Test Cases}\label{tab:author_test_cases}\\
\hline
\textbf{Test Case ID} & \textbf{Test Title} & \textbf{Test Steps} & \textbf{Expected Results} \\
\hline
\endfirsthead
\multicolumn{4}{c}%
{{\bfseries \tablename\ \thetable{} -- continued from previous page}} \\
\hline
\textbf{Test Case ID} & \textbf{Test Title} & \textbf{Test Steps} & \textbf{Expected Results} \\
\hline
\endhead
\hline \multicolumn{4}{|r|}{{Continued on next page}} \\ \hline
\endfoot
\hline
\endlastfoot

\multicolumn{4}{|l|}{\textbf{Sub-Section: Author CRUD Operations}} \\ \hline
TC\_AUTHOR\_CREATE\_001 & Create a new author successfully &
    \begin{itemize}[noitemsep, topsep=0pt, partopsep=0pt, leftmargin=*]
        \item Send a POST request with a valid first and last name.
    \end{itemize} &
    \begin{itemize}[noitemsep, topsep=0pt, partopsep=0pt, leftmargin=*]
        \item A new author is created.
        \item The new author's data is returned.
        \item API returns status 201.
    \end{itemize} \\ \hline
TC\_AUTHOR\_READ\_001 & Retrieve all authors &
    \begin{itemize}[noitemsep, topsep=0pt, partopsep=0pt, leftmargin=*]
        \item Send a GET request to the all authors endpoint.
    \end{itemize} &
    \begin{itemize}[noitemsep, topsep=0pt, partopsep=0pt, leftmargin=*]
        \item A list of all author objects is returned.
        \item API returns status 200.
    \end{itemize} \\ \hline
TC\_AUTHOR\_READ\_002 & Retrieve a single author by valid ID &
    \begin{itemize}[noitemsep, topsep=0pt, partopsep=0pt, leftmargin=*]
        \item Send a GET request with a valid author ID.
    \end{itemize} &
    \begin{itemize}[noitemsep, topsep=0pt, partopsep=0pt, leftmargin=*]
        \item The corresponding author's data is returned.
        \item API returns status 200.
    \end{itemize} \\ \hline
TC\_AUTHOR\_READ\_003 & Attempt to retrieve a non-existent author &
    \begin{itemize}[noitemsep, topsep=0pt, partopsep=0pt, leftmargin=*]
        \item Send a GET request with an invalid author ID.
    \end{itemize} &
    \begin{itemize}[noitemsep, topsep=0pt, partopsep=0pt, leftmargin=*]
        \item Request fails.
        \item An error message "Author not found" is returned.
        \item API returns status 404.
    \end{itemize} \\ \hline
TC\_AUTHOR\_UPDATE\_001 & Update an existing author successfully &
    \begin{itemize}[noitemsep, topsep=0pt, partopsep=0pt, leftmargin=*]
        \item Send a PUT request with a valid author ID and updated data.
    \end{itemize} &
    \begin{itemize}[noitemsep, topsep=0pt, partopsep=0pt, leftmargin=*]
        \item The author's record is updated.
        \item The updated author data is returned.
        \item API returns status 200.
    \end{itemize} \\ \hline
TC\_AUTHOR\_DELETE\_001 & Delete an existing author successfully &
    \begin{itemize}[noitemsep, topsep=0pt, partopsep=0pt, leftmargin=*]
        \item Send a DELETE request with a valid author ID.
    \end{itemize} &
    \begin{itemize}[noitemsep, topsep=0pt, partopsep=0pt, leftmargin=*]
        \item The author's record is deleted.
        \item A success message is returned.
        \item API returns status 200.
    \end{itemize} \\ \hline
\end{longtable}

\newpage
\subsection{Book Management}
\begin{longtable}{|>{\raggedright\arraybackslash}p{4cm}|>{\raggedright\arraybackslash}p{5cm}|>{\raggedright\arraybackslash}p{6cm}|>{\raggedright\arraybackslash}p{6cm}|}
\caption{Book Management Test Cases}\label{tab:book_test_cases}\\
\hline
\textbf{Test Case ID} & \textbf{Test Title} & \textbf{Test Steps} & \textbf{Expected Results} \\
\hline
\endfirsthead
\multicolumn{4}{c}%
{{\bfseries \tablename\ \thetable{} -- continued from previous page}} \\
\hline
\textbf{Test Case ID} & \textbf{Test Title} & \textbf{Test Steps} & \textbf{Expected Results} \\
\hline
\endhead
\hline \multicolumn{4}{|r|}{{Continued on next page}} \\ \hline
\endfoot
\hline
\endlastfoot

\multicolumn{4}{|l|}{\textbf{Sub-Section: Book CRUD Operations}} \\ \hline
TC\_BOOK\_CREATE\_001 & Create a new book successfully &
    \begin{itemize}[noitemsep, topsep=0pt, partopsep=0pt, leftmargin=*]
        \item Send a POST request with all required, valid book details.
    \end{itemize} &
    \begin{itemize}[noitemsep, topsep=0pt, partopsep=0pt, leftmargin=*]
        \item A new book record is created.
        \item The new book data is returned.
        \item API returns status 201.
    \end{itemize} \\ \hline
TC\_BOOK\_READ\_001 & Retrieve all books &
    \begin{itemize}[noitemsep, topsep=0pt, partopsep=0pt, leftmargin=*]
        \item Send a GET request to the all books endpoint.
    \end{itemize} &
    \begin{itemize}[noitemsep, topsep=0pt, partopsep=0pt, leftmargin=*]
        \item A list of all book objects is returned.
        \item API returns status 200.
    \end{itemize} \\ \hline
TC\_BOOK\_UPDATE\_001 & Update an existing book successfully &
    \begin{itemize}[noitemsep, topsep=0pt, partopsep=0pt, leftmargin=*]
        \item Send a PUT request with a valid book ID and updated data.
    \end{itemize} &
    \begin{itemize}[noitemsep, topsep=0pt, partopsep=0pt, leftmargin=*]
        \item The book's record is updated.
        \item The updated book data is returned.
        \item API returns status 200.
    \end{itemize} \\ \hline
TC\_BOOK\_DELETE\_001 & Delete an existing book successfully &
    \begin{itemize}[noitemsep, topsep=0pt, partopsep=0pt, leftmargin=*]
        \item Send a DELETE request with a valid book ID.
    \end{itemize} &
    \begin{itemize}[noitemsep, topsep=0pt, partopsep=0pt, leftmargin=*]
        \item The book's record is deleted.
        \item A success message is returned.
        \item API returns status 200.
    \end{itemize} \\ \hline
\end{longtable}

\newpage
\subsection{Borrowal Management}
\begin{longtable}{|>{\raggedright\arraybackslash}p{4cm}|>{\raggedright\arraybackslash}p{5cm}|>{\raggedright\arraybackslash}p{6cm}|>{\raggedright\arraybackslash}p{6cm}|}
\caption{Borrowal Management Test Cases}\label{tab:borrowal_test_cases}\\
\hline
\textbf{Test Case ID} & \textbf{Test Title} & \textbf{Test Steps} & \textbf{Expected Results} \\
\hline
\endfirsthead
\multicolumn{4}{c}%
{{\bfseries \tablename\ \thetable{} -- continued from previous page}} \\
\hline
\textbf{Test Case ID} & \textbf{Test Title} & \textbf{Test Steps} & \textbf{Expected Results} \\
\hline
\endhead
\hline \multicolumn{4}{|r|}{{Continued on next page}} \\ \hline
\endfoot
\hline
\endlastfoot

\multicolumn{4}{|l|}{\textbf{Sub-Section: Borrowal Operations}} \\ \hline
TC\_BORROW\_CREATE\_001 & Create new borrowal for available book &
    \begin{itemize}[noitemsep, topsep=0pt, partopsep=0pt, leftmargin=*]
        \item Send a POST request with a valid book ID (for an available book) and member ID.
    \end{itemize} &
    \begin{itemize}[noitemsep, topsep=0pt, partopsep=0pt, leftmargin=*]
        \item A new borrowal record is created.
        \item The book's status is updated to 'borrowed'.
        \item API returns status 201.
    \end{itemize} \\ \hline
TC\_BORROW\_CREATE\_002 & Attempt borrowal for unavailable book &
    \begin{itemize}[noitemsep, topsep=0pt, partopsep=0pt, leftmargin=*]
        \item Send a POST request for a book that is already borrowed.
    \end{itemize} &
    \begin{itemize}[noitemsep, topsep=0pt, partopsep=0pt, leftmargin=*]
        \item Request is rejected.
        \item An error message "Book is not available" is returned.
        \item API returns status 400.
    \end{itemize} \\ \hline
TC\_BORROW\_UPDATE\_001 & Update borrowal record (return book) &
    \begin{itemize}[noitemsep, topsep=0pt, partopsep=0pt, leftmargin=*]
        \item Send a PUT request with a valid borrowal ID and a return date.
    \end{itemize} &
    \begin{itemize}[noitemsep, topsep=0pt, partopsep=0pt, leftmargin=*]
        \item The borrowal record is updated.
        \item The associated book's status is updated to 'available'.
        \item API returns status 200.
    \end{itemize} \\ \hline
\end{longtable}

\newpage
\section{Client-Side Functional Tests}

\subsection{Login Form Component}

\begin{longtable}{|>{\raggedright\arraybackslash}p{4cm}|>{\raggedright\arraybackslash}p{5cm}|>{\raggedright\arraybackslash}p{6cm}|>{\raggedright\arraybackslash}p{6cm}|}
\caption{Login Form Test Cases}\label{tab:login_form_test_cases}\\
\hline
\textbf{Test Case ID} & \textbf{Test Title} & \textbf{Test Steps} & \textbf{Expected Results} \\
\hline
\endfirsthead
\multicolumn{4}{c}%
{{\bfseries \tablename\ \thetable{} -- continued from previous page}} \\
\hline
\textbf{Test Case ID} & \textbf{Test Title} & \textbf{Test Steps} & \textbf{Expected Results} \\
\hline
\endhead
\hline \multicolumn{4}{|r|}{{Continued on next page}} \\ \hline
\endfoot
\hline
\endlastfoot

\multicolumn{4}{|l|}{\textbf{Sub-Section: Form Display and Submission}} \\ \hline
TC\_LOGINFORM\_UI\_001 & Verify login form renders correctly &
    \begin{itemize}[noitemsep, topsep=0pt, partopsep=0pt, leftmargin=*]
        \item Navigate to the login page.
    \end{itemize} &
    \begin{itemize}[noitemsep, topsep=0pt, partopsep=0pt, leftmargin=*]
        \item The login form is displayed.
        \item Email and password input fields are visible.
        \item A 'Login' button is visible.
    \end{itemize} \\ \hline
TC\_LOGINFORM\_SUBMIT\_001 & Successful login submission &
    \begin{itemize}[noitemsep, topsep=0pt, partopsep=0pt, leftmargin=*]
        \item Enter a valid email and password into the form fields.
        \item Click the 'Login' button.
    \end{itemize} &
    \begin{itemize}[noitemsep, topsep=0pt, partopsep=0pt, leftmargin=*]
        \item A login request is sent to the backend.
        \item The user is redirected to the dashboard upon success.
    \end{itemize} \\ \hline

\multicolumn{4}{|l|}{\textbf{Sub-Section: Form Validation}} \\ \hline
TC\_LOGINFORM\_VAL\_001 & Attempt submission with invalid email format &
    \begin{itemize}[noitemsep, topsep=0pt, partopsep=0pt, leftmargin=*]
        \item Enter an invalidly formatted email (e.g., "test@test").
        \item Attempt submission.
    \end{itemize} &
    \begin{itemize}[noitemsep, topsep=0pt, partopsep=0pt, leftmargin=*]
        \item Form submission is prevented.
        \item An error message "Email must be a valid email address" is displayed near the email field.
    \end{itemize} \\ \hline
TC\_LOGINFORM\_VAL\_002 & Attempt submission with empty password &
    \begin{itemize}[noitemsep, topsep=0pt, partopsep=0pt, leftmargin=*]
        \item Enter a valid email.
        \item Leave the password field empty.
        \item Attempt to submit the form.
    \end{itemize} &
    \begin{itemize}[noitemsep, topsep=0pt, partopsep=0pt, leftmargin=*]
        \item Form submission is prevented.
        \item An error message "Password is required" is displayed near the password field.
    \end{itemize} \\ \hline
\end{longtable}

\subsection{Data Formatting Utilities}
\begin{longtable}{|>{\raggedright\arraybackslash}p{4cm}|>{\raggedright\arraybackslash}p{5cm}|>{\raggedright\arraybackslash}p{6cm}|>{\raggedright\arraybackslash}p{6cm}|}
\caption{Data Formatting Utility Test Cases}\label{tab:format_util_test_cases}\\
\hline
\textbf{Test Case ID} & \textbf{Test Title} & \textbf{Test Steps} & \textbf{Expected Results} \\
\hline
\endfirsthead
\multicolumn{4}{c}%
{{\bfseries \tablename\ \thetable{} -- continued from previous page}} \\
\hline
\textbf{Test Case ID} & \textbf{Test Title} & \textbf{Test Steps} & \textbf{Expected Results} \\
\hline
\endhead
\hline \multicolumn{4}{|r|}{{Continued on next page}} \\ \hline
\endfoot
\hline
\endlastfoot

\multicolumn{4}{|l|}{\textbf{Sub-Section: Number Formatting}} \\ \hline
TC\_FORMAT\_NUM\_001 & Format number as currency &
    \begin{itemize}[noitemsep, topsep=0pt, partopsep=0pt, leftmargin=*]
        \item Call the `fCurrency` function with a number (e.g., 12345).
    \end{itemize} &
    \begin{itemize}[noitemsep, topsep=0pt, partopsep=0pt, leftmargin=*]
        \item The number is returned as a formatted currency string (e.g., "\$12,345.00").
    \end{itemize} \\ \hline
TC\_FORMAT\_NUM\_002 & Shorten a large number &
    \begin{itemize}[noitemsep, topsep=0pt, partopsep=0pt, leftmargin=*]
        \item Call the `fShortenNumber` function with a large number (e.g., 1234567).
    \end{itemize} &
    \begin{itemize}[noitemsep, topsep=0pt, partopsep=0pt, leftmargin=*]
        \item The number is returned in a compact, human-readable format (e.g., "1.23m").
    \end{itemize} \\ \hline
TC\_FORMAT\_NUM\_003 & Format number as data size &
    \begin{itemize}[noitemsep, topsep=0pt, partopsep=0pt, leftmargin=*]
        \item Call the `fData` function with a number representing bytes (e.g., 1234567).
    \end{itemize} &
    \begin{itemize}[noitemsep, topsep=0pt, partopsep=0pt, leftmargin=*]
        \item The number is converted to a human-readable data size string (e.g., "1.18 MB").
    \end{itemize} \\ \hline
\multicolumn{4}{|l|}{\textbf{Sub-Section: Time Formatting}} \\ \hline
TC\_FORMAT\_TIME\_001 & Format a date string &
    \begin{itemize}[noitemsep, topsep=0pt, partopsep=0pt, leftmargin=*]
        \item Call the `fDate` function with a valid date object or string.
    \end{itemize} &
    \begin{itemize}[noitemsep, topsep=0pt, partopsep=0pt, leftmargin=*]
        \item The date is returned in the format "dd/MM/yyyy".
    \end{itemize} \\ \hline
TC\_FORMAT\_TIME\_002 & Format date as relative time &
    \begin{itemize}[noitemsep, topsep=0pt, partopsep=0pt, leftmargin=*]
        \item Call the `fToNow` function with a date from the past.
    \end{itemize} &
    \begin{itemize}[noitemsep, topsep=0pt, partopsep=0pt, leftmargin=*]
        \item A human-readable relative time string is returned (e.g., "about 2 hours ago").
    \end{itemize} \\ \hline
\end{longtable}

\end{document}
